\section{Pregunta N$^{\circ}$17\qquad Brian Huaman Garcia}

\begin{frame}
	\begin{enumerate}\setcounter{enumi}{16}
		\item

		      Para cada $n\in\mathbb{N}$ se definen
		      \begin{equation*}
			      A_{n}=
			      \begin{pmatrix}
				      1 & 2               \\
				      2 & 4+\frac{1}{n^2}
			      \end{pmatrix}
			      \text{ y }
			      b_{n}=
			      \begin{pmatrix}
				      1 \\
				      2-\frac{1}{n^2}
			      \end{pmatrix}
		      \end{equation*}
		      Se desea resolver el sistema $A_{n}x=b_{n}$.
		      Un estudiante obtiene como resultado
		      \begin{math}
			      x^{\ast}=
				      {\left(10\right)}^{T}
		      \end{math}.

		      \begin{enumerate}[a)]
			      \item

			            Se define el vector residual asociado a esta
			            solución como $r_{n}=A_{n}x^{\ast}-b_{n}$.
			            Calcule $r_{n}$.
			            ¿Podemos decir que para $n$ grande, la solución es
			            razonablemente confiable?
			            Justificar.

			      \item

			            Resolver $A_{n}x=b_{n}$ en forma exacta.
			            Denote por $\widehat{x}_{n}$ la solución exacta y
			            calcule el error relativo de la aproximación
			            $x^{\ast}$.

			      \item

			            Lo razonablemente esperado es que
			            $\widehat{x}_{n}$ converja a $x^{\ast}$ cuando
			            $n$ tiende a infinito.
			            Para este ejemplo, ¿se cumple dicha afirmación?
			            Calcule $K_{\infty}\left(A_{n}\right)$ y explique
			            el resultado obtenido.
		      \end{enumerate}
	\end{enumerate}
	\begin{solution}

		\begin{enumerate}[a)]
			\item

			      .

			\item

			      .

			\item

			      .
		\end{enumerate}
	\end{solution}
\end{frame}
